\documentclass[12pt]{article}
\usepackage{threeparttable}
\usepackage{booktabs}
\usepackage{inputenc, amsmath,amssymb, amsthm}
\usepackage[font=small,labelfont=bf,
   justification=justified,
   format=plain]{caption}
\usepackage{subcaption}
\usepackage{longtable}
\usepackage{adjustbox}
\usepackage{graphicx}
\newtheorem{theorem}{Theorem}
\usepackage{xcolor}
\usepackage{float}
\usepackage[colorlinks=true,
            linkcolor=blue,      % color of internal links (sections, pages, etc)
            citecolor=blue,      % color of citation links
            urlcolor=blue,       % color of URL links
            filecolor=blue]{hyperref}
\usepackage{url}
\usepackage{array}
\usepackage{url}
\usepackage{tikz}
\usetikzlibrary{shapes.geometric, arrows, positioning, fit}
\usetikzlibrary{calc}
\usepackage{titling}
\usepackage{lipsum}
\usepackage{natbib}
\usepackage{setspace}
\usepackage[margin=.75in]{geometry}
\usepackage{xcolor}
\usepackage{amsmath, amssymb, amsfonts}  % loads all AMS symbols & environments

\theoremstyle{definition}
\newtheorem{definition}{Definition}[section]
\usepackage[english]{babel}

\onehalfspacing


\graphicspath{{../../../../../Dropbox/Lawsuits Asymmetry/output}}


\title{Entitlement Pipeline: Data Cleaning}
\author{Tyler Jacobson}
\begin{document}
\maketitle

\section{Goal and Plan}
The goal of this project at this point is to map the full housing construction pipeline. How many units are proposed under the entitlements system, how many get approved and how many get built. The crucial elements of this are:
\begin{enumerate}
    \item Having a project level file. There are multiple entitlements for the same project. We need to be confident that we have a single row for each project. How should we handle an entitlement that was proposed, not built and then they try again. If it was at least five years later or is a different kind of project. 
\end{enumerate}

Key Variables We Need:
\begin{itemize}
    \item Date of proposal, completion 
    \item Location 
    \item What they were asking for: units, free text field, type of entitlement
    \item Level of opposition: appeals, environmental review
\end{itemize}

\section{A Project Level Entitlements File}
We want there to be a single row for every proposed entitlement project in the data. Some questions:
\begin{enumerate}
    \item How can we link across cases? There are APNs and addresses on each project that allow us to link across cases. If we make some mistakes, it seems okay as long as we don't lead to missing merges because of the different APNs. 
    \item What if there are multiple APNs for the same project? This is the same issue as above. 
    \item What if there are multiple distinct projects per APN? Like someone tries something now and they also had tried five years ago. I would call that a failure and then a success, so I would want to have like five year buckets.
    \item How can we validate this data? 
\end{enumerate}

\subsection{Tasks}
\begin{enumerate}
    \item Clean Post 2015 File 
    \item Clean Pre 2015 File and assess how data quality is different across the two
    \item Handle duplicates
    \item Classify type of entitlements
    \item Summary statistics
\end{enumerate}


\subsection{Data Sources}
LA entitlements provides a handful of files. In the pre and post 2015 data, each row is a case. This means that there are environmental, and appeals cases that are tied to a main entitlement case all combined. The first step in cleaning the data is to separate out appeals and environmental cases to be wide on the main case. This lets us observe the full length of time for this cases but 

\subsection{Going from Cases to Applications}
Okay so at this point our data is wide where each row is a case and we observe the environmental review case and the appeals case wide. We know that there are many such cases within a single application. How should we aggregate information?
\begin{itemize}
    \item We would want to know if any of the cases are denied
    \item We would want to know the full length associated with the application
    \item We would want the total hold days across the different cases 
    \item We would want to be sure there was no geographic differences across the cases. We would want to know what all the entitlements asked for are. 
\end{itemize}

\begin{figure}[htbp!]
   \centering
   \includegraphics[width=0.6\textwidth]{entitlements_sum_stats/single_case_application_ts.pdf}
   \caption{Fraction of Applications with Single Case}
  \end{figure}

  \subsection{Units Proposed vs Built}
  Here is an example of a case that has 0 approved units in the data but appears to be approved in my data. 
  \begin{quote}
    DIR-2013-3248-CLQ-ACI
  \end{quote}
  In ZIMAS, this building was built in 2018 with 135 units so it appears indeed that it was completed. This implies that the case denied or withdrawn is probably more reliable.

\subsection{Classifying the Type of Entitlement}
Different entitlements have different levels of heavy lift. 

\section{Understanding the Data}
\begin{itemize}
    \item Planning projects are distributed across the three geo teams with an expedited review section that handles eligible projects
    \item Different suffixes will map to different decisionmakers so the final decisionmaker may be different 
    \item Steps of review: entitlement package review, design review, environmental analysis, applicant submitted findings, recommendation to approve or deny, potential conditions
    \item Design review may only be applicable to certain kinds of projects 
    \item Planner determines whether an exemption is applicable
    \item Initial study determines negative vs mnd vs EIR
    \item Social or economic change cannot be determined to be 
    \item Appendix G has thresholds for CEQA 
    \item Planners seem to do the CEQA reviews
\end{itemize}


\section{Entitlement Timelines}
One case is CPC-2015-4184-GPA-ZC-BL-SPR. It's described as, ``70 CONDOMINIUM UNITS WITHIN NINE SEPARATE BUILDINGS, 4 STORIES HAVING 3 LEVELS OF DWELLINGS OVER ONE AT GRADE PARKING GARAGE.'' The project is asking for a zoning change and a GPA and a SPR.

It was submitted on 11/16/2015 and was accepted one month later. It went to a hearing on 8/30/2016. CPC took action on 12/21/2016. The data says it was assigned on 9/1/2017, which is after the final approval. 

The other case associated with this application is a VTT-73939-CN. This was filed on the same date, accepted on 5/20/2016, assigned on 11/13/2015 and approved on 9/21/2016. This is a vesting tentative tract application (because this project is a subdivision), which locks in zoning rules. 

ENV-2015-4183-MND is the environmental case, which was filed on 11/16/2015 and approved on 8/10/2016. It got assigned on 5/23/16 to William Hughen. 

In one of the PDF's, it's clear that the project was extended. Laura Steele (planner) was listed on the form that was dated on November 19 2018.

Another form says they were granted an extension until 2026 because of a delay in financing the project. ``please be advised that it has been shown that a delay has been or will be unavoidable due to a delay in financing the project May 2, 2017 ''

So here it's clear that Laura Steele only got involved for the extention and for some reason needed to do CEQA again. So the date she was assigned and when she was assigned was not reliable. 

\begin{figure}[H]
   \centering
   \includegraphics[width=0.6\textwidth]{entitlements_sum_stats/hist_gap_assigned.pdf}
   \caption{Gap Between Assignment and Filing}
  \end{figure}

Let's look at a handful of other projects that have these kinds of gaps.


\section{Linking Entitlements to Building Permits}

\end{document}